\subsection{超立方体粗粒化的规范体积与边界切割：体--界字典的离散等周接口}\label{subsec:fold-hypercube-bulk-boundary-gauge-volume}
本小节在与定理 \ref{thm:fiber-gauge-volume-chain-factorization} 一致的纤维规范群口径下，讨论超立方体微态空间 \(\Omega_n=\{0,1\}^n\) 上的任意满射粗粒化 \(f:\Omega_n\twoheadrightarrow X\)。
核心目标是将纤维规范体积密度与界面切割面积联系起来，并以边界邻接多图的回路圆维（自由秩）给出等价表述。

\begin{definition}[超立方体粗粒化的界面面积与规范体积密度]\label{def:hypercube-coarsegraining-area-gauge-volume}
记 \(Q_n=(\Omega_n,E_n)\) 为 \(n\)-维超立方体图（相邻当且仅当 Hamming 距离为 \(1\)），并取满射 \(f:\Omega_n\twoheadrightarrow X\)。
对 \(x\in X\) 记纤维 \(S_x:=f^{-1}(x)\) 与纤维大小 \(d_x:=\card{S_x}\)。
定义纤维规范群
\[
\mathcal G(f):=\mathcal G(f:\Omega_n\twoheadrightarrow X)
=\prod_{x\in X}\mathrm{Sym}(S_x),
\qquad
\card{\mathcal G(f)}=\prod_{x\in X} d_x!,
\]
以及规范体积密度
\[
g(f):=2^{-n}\log\card{\mathcal G(f)}=2^{-n}\sum_{x\in X}\log(d_x!).
\]
再定义界面面积（切割边数）
\[
A(f):=\card{\{\{u,v\}\in E_n:\ f(u)\neq f(v)\}}.
\]
\end{definition}

\begin{theorem}[Shannon--Stirling 双展开与体--界熵缺口恒等式]\label{thm:hypercube-gauge-volume-shannon-stirling-gap}
在定义 \ref{def:hypercube-coarsegraining-area-gauge-volume} 的设定下，令推前分布 \(p_x:=d_x/2^n\)（\(x\in X\)），并记 Shannon 熵
\[
H(p):=-\sum_{x\in X}p_x\log p_x.
\]
则存在显式余项 \(R(f)\ge 0\) 使得严格恒等式成立：
\[
\boxed{\;
g(f)=n\log 2-1-H(p)
\;+\;\frac1{2^{n+1}}\sum_{x\in X}\log(2\pi d_x)
\;+\;R(f).
\;}
\]
并且余项满足统一上界
\[
0\le R(f)\le \frac{|X|}{12\cdot 2^n}.
\]
因此
\[
\bigl|g(f)-(n\log 2-1-H(p))\bigr|
\le
\frac{|X|}{2^{n+1}}\log(2\pi 2^n)+\frac{|X|}{12\cdot 2^n}.
\]
特别地，若 \(|X|/2^n\to 0\)（当 \(n\to\infty\)），则
\(
g(f)=n\log 2-1-H(p)+o(1)
\)。
\end{theorem}
\begin{proof}
对整数 \(m\ge 1\) 取带余项的 Stirling 展开
\[
\log(m!)=m\log m-m+\tfrac12\log(2\pi m)+\theta_m,
\qquad
0\le \theta_m\le \frac1{12m}.
\]
将其代入 \(g(f)=2^{-n}\sum_x\log(d_x!)\) 并逐项求和得
\[
g(f)
=2^{-n}\sum_x\bigl(d_x\log d_x-d_x\bigr)
\;+\;\frac1{2^{n+1}}\sum_x\log(2\pi d_x)
\;+\;2^{-n}\sum_x\theta_{d_x}.
\]
注意 \(\sum_x d_x=2^n\)，故 \(2^{-n}\sum_x d_x=1\)，并且
\[
2^{-n}\sum_x d_x\log d_x
=\sum_x p_x\log(2^n p_x)
=n\log 2+\sum_x p_x\log p_x
=n\log 2-H(p).
\]
令
\(
R(f):=2^{-n}\sum_x\theta_{d_x}
\)，则 \(R(f)\ge 0\) 且
\[
0\le R(f)\le 2^{-n}\sum_x\frac1{12d_x}\le \frac{|X|}{12\cdot 2^n}.
\]
从而得到所述恒等式与误差界。证毕。
\end{proof}

\begin{proposition}[规范体积密度的 Rényi 型上界（由碰撞矩直接控制）]\label{prop:hypercube-gauge-volume-renyi-upperbound}
在定义 \ref{def:hypercube-coarsegraining-area-gauge-volume} 的设定下，令推前分布 \(p_x:=d_x/2^n\)（\(x\in X\)）。
对任意 \(q>1\) 定义 \(q\)-碰撞矩
\[
S_q(f):=\sum_{x\in X} d_x^q.
\]
则规范体积密度满足上界
\begin{equation}\label{eq:hypercube-gauge-volume-renyi-upperbound}
\boxed{\;
g(f)\ \le\ \frac{1}{q-1}\Bigl(\log S_q(f)-n\log 2\Bigr)\ -\ 1\ +\ \frac{|X|}{2^n}.\;}
\end{equation}
\end{proposition}
\begin{proof}
由基本不等式 \(\log(m!)\le m\log m-m+1\)（\(m\ge 1\)）得
\[
g(f)=2^{-n}\sum_{x\in X}\log(d_x!)
\le 2^{-n}\sum_{x\in X}\bigl(d_x\log d_x-d_x+1\bigr)
=\sum_{x\in X}p_x\log d_x-1+\frac{|X|}{2^n}.
\]
又对 \(q>1\)，由 Jensen 不等式（\(\log\) 凹）得到
\[
\sum_{x\in X}p_x\log d_x
=\frac{1}{q-1}\sum_{x\in X}p_x\log d_x^{\,q-1}
\le \frac{1}{q-1}\log\sum_{x\in X}p_x d_x^{\,q-1}
=\frac{1}{q-1}\log\Bigl(\frac{S_q(f)}{2^n}\Bigr),
\]
代回即得 \eqref{eq:hypercube-gauge-volume-renyi-upperbound}。证毕。
\end{proof}

\begin{proposition}[离散全息斯托克不等式：规范体积密度的切割下界]\label{prop:hypercube-discrete-holographic-stokes-gauge-area}
在定义 \ref{def:hypercube-coarsegraining-area-gauge-volume} 的设定下，恒有
\[
\boxed{\;
g(f)\ \ge\ (n\log 2-1)\ -\ 2\log 2\cdot\frac{A(f)}{2^n}\ +\ \frac{\card{X}}{2^n}.
\;}
\]
\end{proposition}

\begin{proof}
\textbf{（阶乘下界）}
对整数 \(m\ge 1\) 有
\[
\log(m!)=\sum_{k=1}^{m}\log k\ \ge\ \int_{1}^{m}\log t\,\dd t\ =\ m\log m-m+1.
\]
故
\[
\sum_{x\in X}\log(d_x!)
\ge
\sum_{x\in X}(d_x\log d_x-d_x+1)
=
\sum_{x\in X}d_x\log d_x-2^n+\card{X}.
\]

\textbf{（超立方体边缘等周不等式）}
对任意 \(S\subseteq \Omega_n\)，记其边界边集
\(
\partial S:=\{\{u,v\}\in E_n:\ u\in S,\ v\notin S\}
\)。
超立方体的标准边缘等周不等式给出
\[
\card{\partial S}
\ge
\card{S}\log_2\!\Bigl(\frac{2^n}{\card{S}}\Bigr)
=
\card{S}\bigl(n-\log_2\card{S}\bigr).
\]
应用于每个纤维 \(S_x\) 得
\[
\card{\partial S_x}\ge d_x\bigl(n-\log_2 d_x\bigr)
\quad\Longrightarrow\quad
d_x\log_2 d_x\ \ge\ nd_x-\card{\partial S_x}.
\]
对 \(x\in X\) 求和并用 \(\sum_x d_x=2^n\)，得
\[
\sum_{x\in X}d_x\log_2 d_x
\ge
n2^n-\sum_{x\in X}\card{\partial S_x}.
\]
注意每条切割边 \(\{u,v\}\)（满足 \(f(u)\neq f(v)\)）恰被计入 \(\sum_x\card{\partial S_x}\) 两次，故 \(\sum_x\card{\partial S_x}=2A(f)\)。
于是
\[
\sum_{x\in X}d_x\log d_x
=
\log 2\cdot \sum_{x\in X}d_x\log_2 d_x
\ge
\log 2\cdot (n2^n-2A(f)).
\]

\textbf{（合并）}
代回阶乘下界并除以 \(2^n\) 得
\[
g(f)
\ge
2^{-n}\Bigl(\log 2\cdot (n2^n-2A(f))-2^n+\card{X}\Bigr),
\]
即所示不等式。证毕。
\end{proof}

\begin{proposition}[面积律的熵上界与二相刚性]\label{prop:hypercube-entropy-area-rigidity}
在定义 \ref{def:hypercube-coarsegraining-area-gauge-volume} 的设定下，令 \(p_x:=d_x/2^n\) 为纤维分布，并记 Shannon 熵
\[
H(p):=-\sum_{x\in X}p_x\log p_x.
\]
则有：
\begin{enumerate}
  \item \label{it:hypercube-entropy-area} 熵--面积上界
  \[
  \boxed{\;
  H(p)\ \le\ 2\log 2\cdot \frac{A(f)}{2^n}.
  \;}
  \]
  \item \label{it:hypercube-area-collapse} 若 \(A(f)=o(2^n)\)（当 \(n\to\infty\)），则 \(\max_x p_x\to 1\)。
  \item \label{it:hypercube-area-saturation} 记 \(\card{E_n}=n2^{n-1}\)。若 \(A(f)=\card{E_n}-o(2^n)\)，则对均匀随机边 \(\{U,V\}\in E_n\) 有 \(\PP(f(U)=f(V))\to 0\)。
\end{enumerate}
\end{proposition}

\begin{proof}
\textbf{（\ref{it:hypercube-entropy-area}）}
由边缘等周不等式，\(\card{\partial S_x}\ge d_x\log_2(2^n/d_x)=d_x\log_2(1/p_x)\)。
求和得
\[
\sum_{x\in X}\card{\partial S_x}
\ge
2^n\sum_{x\in X}p_x\log_2(1/p_x)
=
2^n H_2(p),
\]
其中 \(H_2(p):=-\sum_x p_x\log_2 p_x\) 为以 \(2\) 为底的熵。
又 \(\sum_x\card{\partial S_x}=2A(f)\)，故 \(H(p)=\log 2\cdot H_2(p)\le 2\log 2\cdot A(f)/2^n\)。

\textbf{（\ref{it:hypercube-area-collapse}）}
设 \(p_{\max}=1-\delta\)。则对任意分布有下界 \(H(p)\ge h(\delta)\)，其中
\(
h(\delta):=-\delta\log\delta-(1-\delta)\log(1-\delta)
\)。
由 \ref{it:hypercube-entropy-area} 若 \(A(f)=o(2^n)\) 则 \(H(p)\to 0\)，从而 \(\delta\to 0\)，即 \(p_{\max}\to 1\)。

\textbf{（\ref{it:hypercube-area-saturation}）}
记内部边数
\[
I(f):=\card{\{\{u,v\}\in E_n:\ f(u)=f(v)\}},
\]
则 \(A(f)=\card{E_n}-I(f)\)。
若 \(A(f)=\card{E_n}-o(2^n)\)，则 \(I(f)=o(2^n)\)。
对均匀随机边 \(\{U,V\}\) 有
\[
\PP(f(U)=f(V))=\frac{I(f)}{\card{E_n}}=\frac{o(2^n)}{n2^{n-1}}\to 0.
\]
证毕。
\end{proof}

\begin{definition}[边界邻接多图]\label{def:hypercube-boundary-multigraph}
在定义 \ref{def:hypercube-coarsegraining-area-gauge-volume} 的设定下，定义边界邻接多图 \(\mathsf B_f\) 为：其顶点集为 \(X\)，并对每条切割边 \(\{u,v\}\in E_n\)（满足 \(f(u)\neq f(v)\)）添加一条连接 \(f(u)\) 与 \(f(v)\) 的无向边（允许重边与自环；本构造中自环不出现）。
于是 \(\card{E(\mathsf B_f)}=A(f)\)。
\end{definition}

\begin{proposition}[回路圆维的欧拉恒等式]\label{prop:hypercube-boundary-multigraph-h1-cdim}
在定义 \ref{def:hypercube-boundary-multigraph} 的设定下，记 \(\mathsf B_f\) 的连通分量数为 \(c(\mathsf B_f)\)。
则其一维同调为自由阿贝尔群，且
\[
H_1(\mathsf B_f;\ZZ)\ \cong\ \ZZ^{\,A(f)-\card{X}+c(\mathsf B_f)}.
\]
因此（按圆维在有限生成交换群上的自由秩口径）
\[
\boxed{\;
\cdim\bigl(H_1(\mathsf B_f;\ZZ)\bigr)=A(f)-\card{X}+c(\mathsf B_f).
\;}
\]
\end{proposition}

\begin{proof}
将 \(\mathsf B_f\) 视作一维 CW 复形。其链复形满足 \(C_1\cong \ZZ^{\card{E}}\)、\(C_0\cong \ZZ^{\card{V}}\)，且
\(
\operatorname{rank}H_1=\card{E}-\card{V}+c
\)。
代入 \(\card{E}=A(f)\)、\(\card{V}=\card{X}\) 即得秩公式，并得同构 \(H_1(\mathsf B_f;\ZZ)\cong \ZZ^{\operatorname{rank}H_1}\)。
最后由 \(\cdim(\ZZ^r)=r\) 得圆维等式。证毕。
\end{proof}

\input{sections/appendix/fold_multiplicity/thm__graph-coarsegraining-cycle-rank-decomposition}

\begin{theorem}[边界回路格的 \texorpdfstring{$\Theta$}{Theta} 函数极限律：圆维主项与树数校正]\label{thm:hypercube-cycle-lattice-theta-asymptotic}
令 \(G=(V,E)\) 为有限连通多重图，固定一组取向并将 \(1\)-链视作 \(\ZZ^{E}\subset\RR^{E}\)。
令 \(B\in\ZZ^{V\times E}\) 为取向关联矩阵，并令 \(\tilde B\in\ZZ^{(|V|-1)\times E}\) 为删去一行后的约化关联矩阵。
定义回路格
\[
\Lambda(G):=\ker(\tilde B)\cap\ZZ^{E},
\qquad
r:=\operatorname{rank}_{\ZZ}\Lambda(G)=|E|-|V|+1.
\]
以 \(\RR^{E}\) 的标准内积诱导 Euclid 范数 \(|\cdot|\)，并定义回路格的 theta 函数
\[
\Theta_G(t):=\sum_{z\in\Lambda(G)}e^{-\pi t|z|^2},
\qquad t>0.
\]
记 \(\tau(G)\) 为 \(G\) 的生成树个数（Kirchhoff 复杂度）。
则存在常数 \(c(G)>0\) 使得当 \(t\downarrow 0\) 时有渐近展开
\[
\boxed{\;
\Theta_G(t)=t^{-r/2}\,\tau(G)^{-1/2}\Bigl(1+O\bigl(e^{-c(G)/t}\bigr)\Bigr).
\;}
\]
从而
\[
\boxed{\;
\log\Theta_G(t)=\frac r2\log\frac1t-\frac12\log\tau(G)+O\bigl(e^{-c(G)/t}\bigr).
\;}
\]
特别地，取 \(G=\mathsf B_f\) 时，\(r=\cdim(H_1(\mathsf B_f;\ZZ))\)（命题 \ref{prop:hypercube-boundary-multigraph-h1-cdim}），故上述自由能主项系数即为边界回路圆维。
\end{theorem}
\begin{proof}
\textbf{（i）回路格协体积与树数）}
选取 \(G\) 的一棵生成树 \(T\subseteq E\)。
记 \(E= T\sqcup C\) 为边集分解，其中 \(|T|=|V|-1\)、\(|C|=r\)。
按该分解将 \(\tilde B\) 分块写成
\(
\tilde B=[\tilde B_T\ \ \tilde B_C]
\)。
由于 \(T\) 为树，\(\tilde B_T\in\ZZ^{(|V|-1)\times(|V|-1)}\) 可逆且 \(\det(\tilde B_T)=\pm 1\)（关联矩阵的树子式为 \(\pm 1\)）。
于是对任意 \(u\in\ZZ^{r}\)，向量
\[
z(u):=\begin{pmatrix}
 -\tilde B_T^{-1}\tilde B_C\,u\\
 u
\end{pmatrix}\in\ZZ^{E}
\]
满足 \(\tilde B z(u)=0\)，并且 \(u\mapsto z(u)\) 给出 \(\Lambda(G)\) 的一组 \(\ZZ\)-基。
令 \(Z\in\ZZ^{E\times r}\) 为该基矩阵，则回路格的协体积平方为
\[
\operatorname{covol}(\Lambda(G))^2=\det(Z^\top Z).
\]
写 \(X:=-\tilde B_T^{-1}\tilde B_C\)，则 \(Z=\binom{X}{I_r}\)，从而
\[
\det(Z^\top Z)=\det(I_r+X^\top X)=\det(I_{|V|-1}+XX^\top),
\]
其中最后一步用 Sylvester 行列式恒等式。
进一步
\[
XX^\top=\tilde B_T^{-1}\tilde B_C\tilde B_C^\top(\tilde B_T^{-1})^\top,
\]
故
\[
\det(I_{|V|-1}+XX^\top)
=\det(\tilde B_T\tilde B_T^\top)^{-1}\det(\tilde B_T\tilde B_T^\top+\tilde B_C\tilde B_C^\top).
\]
而 \(\det(\tilde B_T\tilde B_T^\top)=\det(\tilde B_T)^2=1\)，并且
\(
\tilde B_T\tilde B_T^\top+\tilde B_C\tilde B_C^\top=\tilde B\tilde B^\top
\)
为约化 Laplacian。
因此
\[
\operatorname{covol}(\Lambda(G))^2=\det(\tilde B\tilde B^\top).
\]
由 Kirchhoff 矩阵树定理，\(\det(\tilde B\tilde B^\top)=\tau(G)\)，从而
\(
\operatorname{covol}(\Lambda(G))=\sqrt{\tau(G)}
\)。

\textbf{（ii）Poisson 求和与极限）}
令 \(\Lambda\subset\RR^{r}\) 为任意满秩欧氏格，\(\Lambda^\ast\) 为其对偶格。
对高斯函数 \(\phi_t(x):=e^{-\pi t|x|^2}\) 有 Fourier 变换 \(\widehat\phi_t(\xi)=t^{-r/2}e^{-\pi|\xi|^2/t}\)。
应用 Poisson 求和公式得
\[
\Theta_\Lambda(t)=\sum_{z\in\Lambda}\phi_t(z)
=\frac{1}{\operatorname{covol}(\Lambda)}\sum_{y\in\Lambda^\ast}\widehat\phi_t(y)
=t^{-r/2}\operatorname{covol}(\Lambda)^{-1}\Theta_{\Lambda^\ast}(1/t).
\]
当 \(t\downarrow 0\) 时，\(1/t\uparrow\infty\)，并且
\[
\Theta_{\Lambda^\ast}(1/t)=1+O\bigl(e^{-\pi \lambda_1(\Lambda^\ast)^2/t}\bigr),
\]
其中 \(\lambda_1(\Lambda^\ast)\) 为对偶格最短非零向量长度。
取 \(c(G):=\pi\lambda_1(\Lambda(G)^\ast)^2\) 并代入 \(\operatorname{covol}(\Lambda(G))=\sqrt{\tau(G)}\) 即得所述两式。证毕。
\end{proof}

% Energy-shell lattice counting on affine cycle lattices; bit lower bound.

\begin{theorem}[回路圆维驱动的椭球格点计数：固定离散 Stokes 数据后的能量壳微观态熵]\label{thm:graph-energy-shell-lattice-counting}
令 \(G\) 为有限连通无自环多重图，并固定一组取向，使
\[
C_1(G;\ZZ)\cong \ZZ^{E},
\qquad
C_0(G;\ZZ)\cong \ZZ^{V},
\]
其中 \(E:=|E(G)|\)、\(V:=|V(G)|\)。
记边界算子（入射算子）为 \(\partial:C_1(G;\ZZ)\to C_0(G;\ZZ)\)，并记回路圆维
\[
r:=\operatorname{rank}_{\ZZ}H_1(G;\ZZ)=\operatorname{rank}_{\ZZ}\ker(\partial)=E-V+1.
\]
给定任意 \(\delta\in \mathrm{im}(\partial)\subset C_0(G;\ZZ)\)，任选一解 \(\varphi_0\in C_1(G;\ZZ)\) 使得 \(\partial\varphi_0=\delta\)。
取权重 \(w:E(G)\to(0,\infty)\)，并定义加权能量二次型
\[
\mathcal E_w(\varphi):=\sum_{e\in E(G)} w_e\,\varphi_e^2.
\]
考虑能量壳内的整数 \(1\)-形式计数
\[
N_{\delta,w}(E):=\card{\left\{\varphi\in C_1(G;\ZZ):\ \partial\varphi=\delta,\ \mathcal E_w(\varphi)\le E\right\}}.
\]
则存在正定矩阵 \(A\in\mathrm{Sym}_r^+(\RR)\)（仅依赖于 \(G\) 的回路格与权重 \(w\)，且与 \(\varphi_0\) 的选取无关）使得当 \(E\to\infty\) 时
\[
N_{\delta,w}(E)=\frac{\operatorname{Vol}(B_r)}{\sqrt{\det A}}\,E^{r/2}+O\!\left(E^{(r-1)/2}\right),
\]
其中 \(\operatorname{Vol}(B_r)=\pi^{r/2}/\Gamma(r/2+1)\) 为 \(\RR^r\) 中单位欧氏球体积。
因此在固定 \(\delta\) 后，能量 \(E\) 内微观态熵满足
\[
\log N_{\delta,w}(E)=\frac r2\log E-\frac12\log\det A+O(1),
\]
其主阶系数仅由回路圆维 \(r\) 决定。
\end{theorem}
\begin{proof}
由于 \(\delta\in\mathrm{im}(\partial)\)，解集非空且为仿射格
\[
\{\varphi\in C_1(G;\ZZ):\ \partial\varphi=\delta\}=\varphi_0+\ker(\partial).
\]
取 \(\ker(\partial)\) 的一组 \(\ZZ\)-基并记相应基矩阵为 \(T\in\ZZ^{E\times r}\)，使得 \(\ker(\partial)=T\ZZ^r\)。
于是任意解可写为 \(\varphi=\varphi_0+Tk\)（\(k\in\ZZ^r\)）。
记 \(W:=\mathrm{diag}(w_e)\in\mathrm{Sym}_E^+(\RR)\)，则
\[
\mathcal E_w(\varphi_0+Tk)=(\varphi_0+Tk)^{\mathsf T}W(\varphi_0+Tk)
=k^{\mathsf T}(T^{\mathsf T}WT)k+2b^{\mathsf T}k+c,
\]
其中 \(A:=T^{\mathsf T}WT\in\mathrm{Sym}_r^+(\RR)\)、\(b:=T^{\mathsf T}W\varphi_0\)、\(c:=\varphi_0^{\mathsf T}W\varphi_0\)。
完成平方可将约束 \(\mathcal E_w(\varphi_0+Tk)\le E\) 改写为
\[
(k+\kappa)^{\mathsf T}A(k+\kappa)\le E',
\]
其中 \(\kappa\in\RR^r\) 与 \(E'=E+O(1)\) 由 \(A,b,c\) 决定。
因此 \(N_{\delta,w}(E)\) 等于一个以 \(A\) 定义的椭球与平移格 \(\ZZ^r+\kappa\) 的交点数。
对固定 \(A\) 的椭球族，Lipschitz 原理（体积近似）给出格点数等于椭球体积加上边界量级误差：
\[
N_{\delta,w}(E)=\operatorname{Vol}\{x\in\RR^r:\ x^{\mathsf T}Ax\le E'\}+O\!\left((E')^{(r-1)/2}\right).
\]
而体积可显式计算为
\[
\operatorname{Vol}\{x:\ x^{\mathsf T}Ax\le E'\}=\frac{\operatorname{Vol}(B_r)}{\sqrt{\det A}}\,(E')^{r/2}
=\frac{\operatorname{Vol}(B_r)}{\sqrt{\det A}}\,E^{r/2}+O\!\left(E^{(r-1)/2}\right).
\]
将误差并入即可得到主式。
由于 \(A=T^{\mathsf T}WT\) 仅依赖于回路格 \(\ker(\partial)\) 与权重 \(w\)，与特解 \(\varphi_0\) 无关。
对对数式，将主项取对数并吸收相对误差得到 \(\log N_{\delta,w}(E)=\frac r2\log E-\frac12\log\det A+O(1)\)。证毕。
\end{proof}

\begin{corollary}[能量约束下的全息比特下界：圆维给出的有效自由度]\label{cor:graph-energy-shell-bit-lower-bound}
在定理 \ref{thm:graph-energy-shell-lattice-counting} 的设定下，任意能够区分集合
\[
\left\{\varphi\in C_1(G;\ZZ):\ \partial\varphi=\delta,\ \mathcal E_w(\varphi)\le E\right\}
\]
中全部元素的单射编码，其最小信息量（以比特计）满足
\[
\mathrm{bits}(E)\ \ge\ \log_2 N_{\delta,w}(E)=\frac r2\log_2 E+O(1).
\]
\end{corollary}
\begin{proof}
任意单射编码对 \(N_{\delta,w}(E)\) 个不同对象至少需要 \(\log_2 N_{\delta,w}(E)\) 比特。
再代入定理 \ref{thm:graph-energy-shell-lattice-counting} 的渐近式即得。证毕。
\end{proof}

\endinput



\begin{corollary}[回路格球截断的体积主项与 Gödel 外置对数规模]\label{cor:hypercube-cycle-lattice-counting-godel-log}
在定理 \ref{thm:hypercube-cycle-lattice-theta-asymptotic} 的设定下，对 \(R>0\) 定义球截断
\[
\Lambda(G;R):=\{z\in\Lambda(G):\ |z|\le R\}.
\]
则存在常数 \(C_G>0\) 使得当 \(R\to\infty\) 时有渐近计数
\[
\boxed{\;
|\Lambda(G;R)|
=\frac{\operatorname{vol}(B_r)}{\sqrt{\tau(G)}}\,R^r+O(R^{r-1}).
\;}
\]
其中 \(\operatorname{vol}(B_r)\) 为 \(\RR^{r}\) 中单位球体积。
进一步，取任意 Gödel 外置单射 \(\Gamma:\Lambda(G)\hookrightarrow\NN\)，则
\[
\boxed{\;
\log|\Gamma(\Lambda(G;R))|
=r\log R-\tfrac12\log\tau(G)+O(1).
\;}
\]
\end{corollary}
\begin{proof}
记 \(\Lambda:=\Lambda(G)\subset \RR^{E}\)，其秩为 \(r\) 且 \(\operatorname{covol}(\Lambda)=\sqrt{\tau(G)}\)（定理 \ref{thm:hypercube-cycle-lattice-theta-asymptotic}）。
取 \(\Lambda\) 在其张成子空间上的一个基本平行体 \(P\)（可取由一组格基生成），并记其直径为 \(D\)。
令 \(B(R)\) 为该子空间中的半径 \(R\) 球。
\par
\textbf{（体积夹逼）}
由于 \(\{z+P:z\in\Lambda\}\) 构成无交叠的平移铺砌，有
\[
|\Lambda\cap B(R)|\cdot \operatorname{vol}(P)
\le \operatorname{vol}(B(R+D)).
\]
同理，若 \(z\in\Lambda\cap B(R-D)\)，则 \(z+P\subset B(R)\)，因此
\[
\operatorname{vol}(B(R-D))
\le |\Lambda\cap B(R)|\cdot \operatorname{vol}(P).
\]
合并并用 \(\operatorname{vol}(P)=\operatorname{covol}(\Lambda)\) 得
\[
\frac{\operatorname{vol}(B(R-D))}{\operatorname{covol}(\Lambda)}
\le |\Lambda(G;R)|
\le \frac{\operatorname{vol}(B(R+D))}{\operatorname{covol}(\Lambda)}.
\]
利用 \(\operatorname{vol}(B(R))=\operatorname{vol}(B_r)R^r\) 与二项展开，
\[
\operatorname{vol}(B(R\pm D))
=\operatorname{vol}(B_r)\,(R\pm D)^r
=\operatorname{vol}(B_r)R^r+O(R^{r-1}),
\]
从而得到首式，且隐含常数仅依赖于 \(G\)。
\par
\textbf{（Gödel 对数）}
由于 \(\Gamma\) 为单射，\(|\Gamma(\Lambda(G;R))|=|\Lambda(G;R)|\)。
对首式取对数即得
\[
\log|\Lambda(G;R)|
=\log\Bigl(\frac{\operatorname{vol}(B_r)}{\sqrt{\tau(G)}}R^r\bigl(1+O(R^{-1})\bigr)\Bigr)
=r\log R-\tfrac12\log\tau(G)+O(1).
\]
证毕。
\end{proof}

\begin{corollary}[回路圆维与规范体积的互补下界]\label{cor:hypercube-loop-cdim-gauge-volume-complement}
在命题 \ref{prop:hypercube-discrete-holographic-stokes-gauge-area} 与命题 \ref{prop:hypercube-boundary-multigraph-h1-cdim} 的设定下，有
\[
\boxed{\;
g(f)\ \ge\ (n\log 2-1)\ -\ 2\log 2\cdot\frac{\cdim(H_1(\mathsf B_f;\ZZ))+\card{X}-c(\mathsf B_f)}{2^n}\ +\ \frac{\card{X}}{2^n}.
\;}
\]
\end{corollary}

\begin{proof}
由命题 \ref{prop:hypercube-boundary-multigraph-h1-cdim} 得
\(
A(f)=\cdim(H_1(\mathsf B_f;\ZZ))+\card{X}-c(\mathsf B_f)
\)。
代入命题 \ref{prop:hypercube-discrete-holographic-stokes-gauge-area} 即得。证毕。
\end{proof}

\begin{corollary}[面积密度极限下的线性面积律]\label{cor:hypercube-linear-area-law-density}
取一列满射粗粒化 \(f_n:\Omega_n\twoheadrightarrow X_n\)。
记面积密度
\[
a_n:=\frac{A(f_n)}{n2^{n-1}}\in[0,1],
\]
以及每比特规范体积密度
\[
\bar g_n:=\frac{g(f_n)}{n}.
\]
若 \(\card{X_n}/2^n\to 0\)，则
\[
\boxed{\;
\liminf_{n\to\infty}\bar g_n\ \ge\ (1-\limsup_{n\to\infty}a_n)\log 2.
\;}
\]
特别地，若 \(a_n\to a\in[0,1]\)，则 \(\liminf_{n\to\infty}\bar g_n\ge (1-a)\log 2\)。
\end{corollary}

\begin{proof}
由命题 \ref{prop:hypercube-discrete-holographic-stokes-gauge-area}，
\[
\bar g_n
\ge
\log 2-\frac{1}{n}-2\log 2\cdot\frac{A(f_n)}{n2^n}+\frac{\card{X_n}}{n2^n}.
\]
注意 \(\frac{A(f_n)}{2^n}=\frac{n}{2}a_n\)，故
\[
\bar g_n\ge (1-a_n)\log 2-\frac{1}{n}+\frac{\card{X_n}}{n2^n}.
\]
取 \(\liminf\) 并用 \(\card{X_n}/2^n\to 0\) 即得结论。证毕。
\end{proof}

\paragraph{Gödel--Stokes 体--界字典：通量偏差、Walsh 全息与回路扭量}
\label{par:fold-hypercube-godel-stokes-walsh-hodge}
本段在超立方体粗粒化的边界邻接多图 \(\mathsf B_f\) 上，给出一条完全离散的“边界 \(\to\) 体积”读出链：单坐标方向的有向通量严格等于体积偏差；Fourier--Walsh 系数可由边界有向 \(1\)-形式积分全息读出；固定体积偏差后边界层仍保留的自由度精确形成以 \(H_1(\mathsf B_f;\ZZ)\) 为平移群的扭量；并且上述数据可由素数寄存器作乘法外置，同时在加权能量下诱导显式椭球可行域。

\begin{proposition}[单坐标离散 Stokes 恒等式：有向通量等于体积偏差]\label{prop:fold-hypercube-godel-stokes-flux-bias}
设 \(n\ge 1\) 且 \(\Omega_n:=\{0,1\}^n\)。
对 \(i\in\{1,\dots,n\}\)，记 \(\omega^{(i)}\) 为将 \(\omega\) 的第 \(i\) 位翻转后的顶点。
对任意 \(S\subseteq\Omega_n\)，定义第 \(i\) 坐标方向的有向边界计数
\[
N_i^+(S):=\#\{\omega\in S:\ \omega_i=0,\ \omega^{(i)}\notin S\},
\qquad
N_i^-(S):=\#\{\omega\in S:\ \omega_i=1,\ \omega^{(i)}\notin S\},
\]
并定义体积偏差
\[
b_i(S):=\#(S\cap\{\omega_i=0\})-\#(S\cap\{\omega_i=1\})
=\sum_{\omega\in S}(-1)^{\omega_i}\in\ZZ.
\]
则对一切 \(i\) 有严格恒等式
\[
\boxed{\;N_i^+(S)-N_i^-(S)=b_i(S)\;}
\]
从而无向边界数 \(c_i(S):=N_i^+(S)+N_i^-(S)\) 满足
\[
c_i(S)\ge |b_i(S)|.
\]
\end{proposition}
\begin{proof}
固定 \(i\)。映射 \(\{\omega:\omega_i=0\}\to\{\omega:\omega_i=1\}\)、\(\omega\mapsto\omega^{(i)}\) 为双射，故
\[
N_i^+(S)-N_i^-(S)
=\sum_{\omega:\omega_i=0}\bigl(\mathbf 1_S(\omega)-\mathbf 1_S(\omega^{(i)})\bigr)
=\#(S\cap\{\omega_i=0\})-\#(S\cap\{\omega_i=1\})
=b_i(S).
\]
再由 \(|a+b|\ge ||a|-|b||\)（取 \(a=N_i^+(S),b=N_i^-(S)\)）得 \(c_i(S)\ge |b_i(S)|\)。证毕。
\end{proof}

\begin{corollary}[\texorpdfstring{$\ell^1$}{l1} 边界下界与坐标单调性的等号刻画]\label{cor:fold-hypercube-boundary-l1-bias-monotone}
在命题 \ref{prop:fold-hypercube-godel-stokes-flux-bias} 的记号下，记 \(\partial S\) 为超立方体图上 \(S\) 的无向边界边集（定义如 \eqref{eq:hypercube-gauge-volume-renyi-upperbound} 的证明中所用）。
则
\[
|\partial S|=\sum_{i=1}^n c_i(S)\ \ge\ \sum_{i=1}^n |b_i(S)|.
\]
并且对给定的 \(i\) 有
\[
c_i(S)=|b_i(S)|
\quad\Longleftrightarrow\quad
N_i^+(S)N_i^-(S)=0,
\]
亦即第 \(i\) 坐标方向的全部边界边在同一方向上外流。
若对所有 \(i\) 同时取等号，则存在坐标翻转自同构
\(
\omega_i\mapsto 1-\omega_i
\)
使得变换后的集合成为坐标偏序下的上闭集（逐坐标单调）。
\end{corollary}
\begin{proof}
第一条由逐坐标不等式 \(c_i(S)\ge |b_i(S)|\) 求和得到。
等号条件来自 \(|a+b|=||a|-|b||\iff ab=0\)（\(a,b\ge 0\)）。
当每个坐标方向的边界边均同向时，按坐标逐层检验可得：在适当坐标翻转后，沿任一坐标由 \(0\) 到 \(1\) 的邻接翻转不会离开集合，从而集合为上闭。证毕。
\end{proof}

\begin{proposition}[Gödel 边界通量码与乘法体积读出]\label{prop:fold-hypercube-godel-flux-multiplicative-volume}
取一组互异素数 \((p_i)_{i=1}^n\)，并定义 Gödel 边界通量码
\[
\Phi_p(S):=\prod_{i=1}^n p_i^{\,N_i^+(S)-N_i^-(S)}\in\QQ_{>0}.
\]
则 \(\Phi_p(S)=\prod_{i=1}^n p_i^{\,b_i(S)}\)，因而仅由体积偏差向量 \((b_1(S),\dots,b_n(S))\) 决定。
进一步定义顶点 Gödel 编码
\[
G(\omega):=\prod_{i=1}^n p_i^{\omega_i}\in\NN_{>0},
\]
则有严格恒等式
\[
\boxed{\;
\prod_{\omega\in S}G(\omega)
=\prod_{i=1}^n p_i^{\,\#(S\cap\{\omega_i=1\})}
=\prod_{i=1}^n p_i^{\frac{|S|-b_i(S)}{2}}\in\NN_{>0}.
\;}
\]
因此 \(\prod_{\omega\in S}G(\omega)\) 被 \(|S|\) 与 \(\Phi_p(S)\) 等价确定（在对数坐标上为一条线性关系）。
\end{proposition}
\begin{proof}
由命题 \ref{prop:fold-hypercube-godel-stokes-flux-bias} 得 \(N_i^+-N_i^-=b_i\)，代回即得 \(\Phi_p\) 的表达式。
另一方面，
\[
\prod_{\omega\in S}G(\omega)
=\prod_{\omega\in S}\prod_{i=1}^n p_i^{\omega_i}
=\prod_{i=1}^n p_i^{\sum_{\omega\in S}\omega_i}
=\prod_{i=1}^n p_i^{\#(S\cap\{\omega_i=1\})}.
\]
又由
\(
b_i(S)=\#(S\cap\{\omega_i=0\})-\#(S\cap\{\omega_i=1\})=|S|-2\#(S\cap\{\omega_i=1\})
\)
得 \(\#(S\cap\{\omega_i=1\})=(|S|-b_i(S))/2\)，代回即得结论。证毕。
\end{proof}

\begin{corollary}[粗粒化切割面积的偏差下界]\label{cor:fold-hypercube-cutedges-bias-l1}
沿用定义 \ref{def:hypercube-coarsegraining-area-gauge-volume} 的满射粗粒化 \(f:\Omega_n\twoheadrightarrow X\)，并记纤维 \(S_x=f^{-1}(x)\)。
则
\[
2A(f)=\sum_{x\in X}|\partial S_x|
\ \ge\
\sum_{x\in X}\sum_{i=1}^n |b_i(S_x)|.
\]
\end{corollary}
\begin{proof}
对每个纤维 \(S_x\) 用推论 \ref{cor:fold-hypercube-boundary-l1-bias-monotone} 的下界 \(|\partial S_x|\ge \sum_i |b_i(S_x)|\)。
再注意每条切割边恰在两条纤维边界中各计一次，故 \(\sum_x|\partial S_x|=2A(f)\)。证毕。
\end{proof}

\begin{theorem}[加权 Walsh 模式的广义 Stokes 公式：Fourier 系数的边界全息读出]\label{thm:fold-hypercube-weighted-walsh-stokes}
在超立方体上赋予坐标边权 \(w_i>0\)，并定义加权拉普拉斯算子
\[
(\Delta_w\psi)(\omega):=\sum_{i=1}^n w_i\bigl(\psi(\omega)-\psi(\omega^{(i)})\bigr).
\]
对非空指标集 \(I\subseteq\{1,\dots,n\}\)，令 Walsh 字符
\[
\chi_I(\omega):=(-1)^{\sum_{i\in I}\omega_i}.
\]
则 \(\chi_I\) 为 \(\Delta_w\) 的本征函数，且
\[
\Delta_w\chi_I=2\Bigl(\sum_{i\in I}w_i\Bigr)\chi_I.
\]
对任意 \(S\subseteq\Omega_n\)，令 \(\partial^{\to}S\) 表示从 \(S\) 指向补集的外法向有向边集合。
定义边界上的有向 \(1\)-形式
\[
\alpha_{I,w}(\omega\to\omega^{(i)})
:=\frac{w_i}{2\sum_{j\in I}w_j}\bigl(\chi_I(\omega^{(i)})-\chi_I(\omega)\bigr),
\qquad (\omega\to\omega^{(i)})\in\partial^{\to}S.
\]
则成立广义离散 Stokes 公式
\[
\boxed{\;\sum_{\omega\in S}\chi_I(\omega)=\sum_{e\in \partial^{\to}S}\alpha_{I,w}(e).\;}
\]
特别地，对 \(f=\mathbf 1_S\) 的 Fourier--Walsh 系数
\(
\widehat f(I)=2^{-n}\sum_{\omega\in\Omega_n}f(\omega)\chi_I(\omega)
\)
有边界积分表达
\[
2^n\widehat{\mathbf 1_S}(I)=\sum_{e\in \partial^{\to}S}\alpha_{I,w}(e).
\]
\end{theorem}
\begin{proof}
由本征关系，
\[
2\Bigl(\sum_{i\in I}w_i\Bigr)\sum_{\omega\in S}\chi_I(\omega)
=\sum_{\omega\in S}(\Delta_w\chi_I)(\omega)
=\sum_{\omega\in S}\sum_{i=1}^n w_i\bigl(\chi_I(\omega)-\chi_I(\omega^{(i)})\bigr).
\]
将内层求和按无向边配对，内部边贡献相消，仅剩跨越 \(S\) 与补集的边；并按外向取向把每条割边写成 \((\omega\to\omega^{(i)})\in\partial^{\to}S\)，得到
\[
\sum_{\omega\in S}(\Delta_w\chi_I)(\omega)
=\sum_{(\omega\to\omega^{(i)})\in\partial^{\to}S}w_i\bigl(\chi_I(\omega)-\chi_I(\omega^{(i)})\bigr).
\]
两边同除 \(2\sum_{i\in I}w_i\) 即得所述公式（与 \(\alpha_{I,w}\) 仅差号，已由定义吸收）。证毕。
\end{proof}

\begin{proposition}[Gödel--Hodge 扭量：固定偏差后的边界自由度由回路群承载]\label{prop:fold-hypercube-godel-hodge-torsor}
沿用定义 \ref{def:hypercube-boundary-multigraph} 的边界邻接多图 \(\mathsf B_f\)，并取其一维链群 \(C_1(\mathsf B_f;\ZZ)\) 与边界算子 \(\partial:C_1\to C_0\)。
对每个坐标 \(i\)，定义纤维偏差向量
\[
b^{(i)}\in C_0(\mathsf B_f;\ZZ),
\qquad
b^{(i)}_x:=b_i(S_x).
\]
则存在一个自然整数流 \(F^{(i)}\in C_1(\mathsf B_f;\ZZ)\)（由超立方体的第 \(i\) 坐标割边按 \(0\to 1\) 取向并推前到 \(\mathsf B_f\) 得到），满足
\[
\partial F^{(i)}=b^{(i)}.
\]
并且若解集
\[
\mathcal S_i:=\{F\in C_1(\mathsf B_f;\ZZ):\ \partial F=b^{(i)}\}
\]
非空，则其为仿射格（torsor）
\[
\mathcal S_i=F^{(i)}+\ker\partial,
\qquad
\ker\partial\cong H_1(\mathsf B_f;\ZZ).
\]
从而对全部坐标同时固定 \((b^{(1)},\dots,b^{(n)})\) 时，残留的全局边界自由度构成以 \(H_1(\mathsf B_f;\ZZ)^n\) 为平移群的扭量；其自由秩为
\[
\operatorname{rank}_{\ZZ}H_1(\mathsf B_f;\ZZ)^n
=n\bigl(A(f)-|X|+c(\mathsf B_f)\bigr)
\]
（命题 \ref{prop:hypercube-boundary-multigraph-h1-cdim}）。
\end{proposition}
\begin{proof}
把超立方体的第 \(i\) 坐标割边按 \(0\to 1\) 方向计为有向 \(1\)-链，并通过端点纤维标记 \(f:\Omega_n\to X\) 推前到 \(\mathsf B_f\) 的边链，得 \(F^{(i)}\)。
对任一顶点 \(x\in X\)，\((\partial F^{(i)})_x\) 等于从纤维 \(S_x\) 发出的 \(0\to 1\) 边数减去 \(1\to 0\) 边数，按命题 \ref{prop:fold-hypercube-godel-stokes-flux-bias} 恰为 \(b_i(S_x)=b_x^{(i)}\)。
若 \(F,F'\) 同为解，则 \(\partial(F-F')=0\)，故差落入 \(\ker\partial\)；反之加上任意 \(Z\in\ker\partial\) 仍为解，得到扭量结构。
最后 \(\ker\partial\cong H_1(\mathsf B_f;\ZZ)\) 与秩公式由命题 \ref{prop:hypercube-boundary-multigraph-h1-cdim} 给出。证毕。
\end{proof}

\begin{proposition}[回路自由度的素数预算下界：自由秩不可压缩]\label{prop:fold-hypercube-h1-prime-budget-minimal}
沿用命题 \ref{prop:fold-hypercube-godel-hodge-torsor} 记号，并令
\[
r:=\operatorname{rank}_{\ZZ}H_1(\mathsf B_f;\ZZ)=A(f)-|X|+c(\mathsf B_f).
\]
则 \(H_1(\mathsf B_f;\ZZ)^n\cong \ZZ^{nr}\)。
取互异素数族 \(\{q_{i,j}\}_{1\le i\le n,\,1\le j\le r}\)，并定义映射
\[
\Gamma:\ \ZZ^{nr}\longrightarrow \QQ_{>0},
\qquad
(a_{i,j})\longmapsto \prod_{i=1}^n\prod_{j=1}^r q_{i,j}^{a_{i,j}}.
\]
则 \(\Gamma\) 为单射群同态。
并且若 \(A\) 为任意自由阿贝尔群且存在单射 \(\ZZ^{nr}\hookrightarrow A\)，则必有 \(\operatorname{rank}_{\ZZ}A\ge nr\)；
等价地，以互异素数作为独立乘法通道对 \(H_1(\mathsf B_f;\ZZ)^n\) 作 Gödel 外置时，所需通道数下界恰为 \(nr\)。
\end{proposition}
\begin{proof}
由唯一分解定理，\(\langle q_{i,j}\rangle\subset\QQ_{>0}\) 的自由部分同构于 \(\ZZ^{nr}\)，故 \(\Gamma\) 单射。
对最小性，任意单射 \(B\hookrightarrow A\)（\(A,B\) 自由阿贝尔）在 \(\QQ\)-张量下诱导线性单射 \(B\otimes\QQ\hookrightarrow A\otimes\QQ\)，故 \(\dim_{\QQ}(B\otimes\QQ)\le \dim_{\QQ}(A\otimes\QQ)\)，即 \(\operatorname{rank}_{\ZZ}A\ge \operatorname{rank}_{\ZZ}B=nr\)。证毕。
\end{proof}

\begin{proposition}[加权边界能量诱导的椭球约束]\label{prop:fold-hypercube-weighted-energy-ellipsoid-bias}
对权重 \(w_i>0\) 定义加权边界能量
\[
E_w(S):=\sum_{i=1}^n w_i\,c_i(S),
\]
其中 \(c_i(S)\) 为第 \(i\) 坐标方向无向边界数（命题 \ref{prop:fold-hypercube-godel-stokes-flux-bias}）。
则成立二次型下界
\[
\boxed{\;
E_w(S)\ \ge\ 2^{1-n}\sum_{i=1}^n w_i\,b_i(S)^2.
\;}
\]
从而在给定 \(E_w(S)\) 时，偏差向量 \(b(S)\) 必落入轴对齐椭球
\[
\sum_{i=1}^n w_i\,b_i^2\ \le\ 2^{n-1}E_w(S).
\]
\end{proposition}
\begin{proof}
令 \(f=\mathbf 1_S\)。由于 \(f(\omega)-f(\omega^{(i)})\in\{0,\pm 1\}\)，有
\[
E_w(S)=\sum_{i=1}^n w_i\sum_{\omega\in\Omega_n}\bigl(f(\omega)-f(\omega^{(i)})\bigr)^2.
\]
将其按 \(2^{-n}\) 归一化得到加权 Dirichlet 形式
\(
\mathcal E_w(f)=2^{-n}E_w(S)
\)。
在 Walsh 基底上（参见定理 \ref{thm:discussion-hypercube-stokes-fourier-binomial} 的谱层展开口径），有
\[
\mathcal E_w(f)=\sum_{\varnothing\ne I\subseteq[n]}2\Bigl(\sum_{i\in I}w_i\Bigr)\,|\widehat f(I)|^2
\ \ge\
\sum_{i=1}^n 2w_i\,|\widehat f(\{i\})|^2.
\]
又 \(\widehat{\mathbf 1_S}(\{i\})=2^{-n}\sum_{\omega\in S}(-1)^{\omega_i}=b_i(S)/2^n\)，故
\[
2^{-n}E_w(S)\ge \sum_{i=1}^n 2w_i\frac{b_i(S)^2}{2^{2n}},
\]
移项即得结论。证毕。
\end{proof}

\endinput



\endinput
